\section{Joint estimates on sparse balls}\label{sec:sparse}

The factor $N^{\gamma_0}$ in Proposition~\ref{prop:discrete} is a
small loss on a bounded region.  To remove it later, we need an estimate
for many distant regions with a single data norm on the right.  The
directions in which two translations fail to oscillate are confined to
a small frequency set.  Discarding that set costs little in $L^\infty$;
the remaining translations are almost orthogonal at the initial packet
scale. This joint estimate supplies the input to the sparse-ball
strategy for removing a spatial loss \cite{Tao1999Epsilon,BourgainGuth2011,Kim2017};
we prove the required estimate for \(L^3\) data with the precise
separation and uniformity below.

\begin{proposition}[Sparse balls]\label{prop:sparse}
There is a constant $\mu_0\in(0,1)$ depending only on the fixed graph
setup with the following property.  Suppose that $h$ is supported in
$Q_0$, $\|h\|_3\leq1$, and $0<\mu\leq\mu_0$.  Let $\mathcal B$ be a
finite family of balls of common radius $R$ such that
\[
 \#\mathcal B\leq\mu^{-8},\qquad R\geq\mu^{-100},\qquad
 |x_B-x_{\widetilde B}|\geq\mu^{-40}R
 \quad(B\ne\widetilde B),
\]
where $x_B$ denotes the center of $B$.  Choose a dyadic integer $N$ with
$4R\leq N^2<16R$.  For $0<\kappa<1/10$ and $0<\eta,s\leq1$, put
$\gamma_0=10\kappa+\eta+3s/2$.  Then
\begin{equation}\label{eq:sparse-estimate}
 \left|\{|F|>\mu\}\cap\bigcup_{B\in\mathcal B}B\right|
 \leq C\mu^{-3}P(\kappa,\eta)M(s)^{3/2}
       N^{\gamma_0}\log^2(2+N).
\end{equation}
The constant $C$ and the threshold $\mu_0$ are independent of
$\kappa,\eta,s$, of the centers, and of their maximum separation.
\end{proposition}

\begin{proof}
Write $\delta=N^{-1}$ and $D=D_{\mathrm f}$.  Every window
$\chi_\delta^\nu$ meeting $Q_0$ lies in a fixed compact neighborhood
$Q_1$ of $Q_0$, since $\delta\leq1$.  Fix
\[
 G=\sup_{Q_1}|\nabla\phi|,\qquad Z>4(2+G).
\]
The initial spatial labels will satisfy $|z|\leq Z$ and the terminal
spatial labels $|z'|\leq1$.  Thus all label bounds used in
Proposition~\ref{prop:discrete} are fixed independently of the small
exponents.  Throughout the proof we decrease $\mu_0$, if necessary,
only in terms of this fixed setup.  In particular, we may require $N$
to exceed any fixed lower threshold needed below.  The hypotheses give
\begin{equation}\label{eq:sparse-scale}
 \delta\leq\tfrac12\mu^{50},\qquad
 \mu^{-8}\leq N,\qquad
 H\geq\mu^{-40}R>\tfrac1{16}\mu^{-40}N^2
\end{equation}
whenever $H$ is the separation of two distinct centers.

\emph{Removing resonant frequencies.}
Set $a_B=x_B-(0,0,N^2/2)$.  For each pair of distinct balls write
$a_B-a_{\widetilde B}=(q,s_0)\in\mathbb R^2\times\mathbb R$ and
$H=|(q,s_0)|$.  Its resonant set is
\[
 \Omega_{B,\widetilde B}
 =\{\xi\in Q_0:|q+s_0\nabla\phi(\xi)|\leq H\mu^{20}\}.
\]
If this set is nonempty, then
\[
 H\leq |q|+|s_0|
 \leq H\mu^{20}+(G+1)|s_0|,
\]
so $|s_0|\geq H/[2(G+1)]$ for sufficiently small $\mu_0$.
For two points of $\Omega_{B,\widetilde B}$,
Equation~\eqref{eq:velocity-monotonicity} gives
\[
 c|\xi-\zeta|
 \leq 2H\mu^{20}/|s_0|\leq C\mu^{20}.
\]
Consequently each resonant set has area at most $C\mu^{40}$.
The union $\Omega$ over all pairs therefore has area at most
$C\mu^{24}$.  Define $h_*=h\mathbf1_{Q_0\setminus\Omega}$, and let
$F_*$ be the extension defining $F$ with $h$ replaced by $h_*$.  By
H\"older's inequality,
\begin{equation}\label{eq:sparse-excision}
 \|F-F_*\|_\infty
 \leq \|\chi_1^{\nu_*}\|_\infty\,\|h\|_3|\Omega|^{2/3}
 \leq C\mu^{16}\leq\mu/2.
\end{equation}
The powers $20$, $40$, and $100$ provide ample room both for this
discarded mass and for the phase estimates below; their precise values
are not essential.

\emph{A joint bound for the initial coefficients.}
Let
\[
 \mathcal U=[0,D^{-1})^2\times[0,1),\qquad
 \Gamma_\nu=\operatorname{supp}\chi_\delta^\nu.
\]
For $u\in\mathcal U$ and $B\in\mathcal B$ define
\begin{equation}\label{eq:sparse-analysis}
 \begin{split}
 h_{B,u}(\xi)
   &=h_*(\xi)e\big((a_B+u)\cdot(\xi,\phi(\xi))\big),\\
 d^{B,u}_{\nu,z}
   &=\int h_{B,u}(\xi)\chi_\delta^\nu(\xi)
                         e(N^2z\cdot\xi)\,d\xi,
       \qquad z\in D^{-1}\delta\mathbb Z^2,\\
 b_\nu^{B,u}&=\sum_{|z|\leq Z}|d^{B,u}_{\nu,z}|^2.
 \end{split}
\end{equation}
Only angular windows meeting $Q_0$ need be retained.  We claim that,
uniformly in $u$,
\begin{equation}\label{eq:sparse-joint-mass}
 \sum_{B\in\mathcal B}b_\nu^{B,u}
 \leq C\delta^2\int_{\Gamma_\nu}|h_*|^2.
\end{equation}

If $h_*\chi_\delta^\nu=0$ almost everywhere, the claim is immediate.
Otherwise choose a point $\xi_g\in\Gamma_\nu\cap(Q_0\setminus\Omega)$.
For this $\nu$, define the operator on \emph{unrestricted}
$L^2(\mathbb R^2)$ by
\[
 (T_\nu f)_{B,z}
 =\int f(\xi)\chi_\delta^\nu(\xi)
   e\big((a_B+u)\cdot(\xi,\phi(\xi))+N^2z\cdot\xi\big)\,d\xi,
 \qquad |z|\leq Z.
\]
We will apply it to $f=h_*\mathbf1_{\Gamma_\nu}$.  This enlargement
of the operator domain is useful: the kernel of $T_\nu T_\nu^*$ has
the smooth amplitude $(\chi_\delta^\nu)^2$, with no indicator of the
excised set.

For a single ball block, Parseval's identity on the square of side
$D\delta$ containing $\Gamma_\nu$ gives
\[
 \sum_{z\in D^{-1}\delta\mathbb Z^2}|(T_\nu f)_{B,z}|^2
 = (D\delta)^2\int |f(\xi)\chi_\delta^\nu(\xi)|^2\,d\xi
 \leq C\delta^2\|f\|_2^2.
\]
Here the full spatial sum is understood before its restriction to
$|z|\leq Z$; its Fourier frequencies are
$N^2z\in(D\delta)^{-1}\mathbb Z^2$.  It follows that each diagonal
ball block of $T_\nu T_\nu^*$ has operator norm at most $C\delta^2$.

For distinct $B,\widetilde B$, an entry of this Gram operator is
\[
 \int(\chi_\delta^\nu(\xi))^2
 e\big((q+N^2(z-\widetilde z))\cdot\xi+s_0\phi(\xi)\big)\,d\xi.
\]
To estimate it, put $\xi=\nu+\delta v$ and denote the resulting
phase by $\Psi(v)$.  Let
\[
 v_0=\frac{q+s_0\nabla\phi(\xi_g)}
              {|q+s_0\nabla\phi(\xi_g)|},\qquad
 \Lambda=\delta H\mu^{20}.
\]
The direction $v_0$ is defined because $\xi_g$ avoids every resonant
set.  On the entire support of the smooth amplitude $\chi(v)^2$,
\begin{align}
 D_{v_0}\Psi(v)
 &\geq \delta\big(H\mu^{20}-C H\delta-2ZN^2\big)
 \geq \Lambda/2,                                      \label{eq:sparse-phase-first}\\
 |\partial^\alpha\Psi(v)|
 &\leq C_\alpha H\delta^2\leq C_\alpha\Lambda
 \quad (|\alpha|\geq2),\qquad
 \Lambda\geq\tfrac1{16}N\mu^{-20}.                    \label{eq:sparse-phase-higher}
\end{align}
For Equation~\eqref{eq:sparse-phase-first}, the two errors relative
to $H\mu^{20}$ are at most $C\mu^{30}$ and $32Z\mu^{20}$ by
Equation~\eqref{eq:sparse-scale}; their sum is at most $1/2$ after
fixing $\mu_0$.  Equation~\eqref{eq:sparse-phase-higher} follows
from the bounded higher derivatives of $\phi$ and
$\delta\mu^{-20}\leq\mu^{30}/2$.

Repeated integration by parts in the fixed direction $v_0$ now gives,
for each fixed integer $k\geq1$,
\begin{equation}\label{eq:sparse-kernel-decay}
 \left|\int(\chi_\delta^\nu)^2
   e\big((q+N^2(z-\widetilde z))\cdot\xi+s_0\phi(\xi)\big)
       \,d\xi\right|
 \leq C_k\delta^2\Lambda^{-k}
 \leq C_k\delta^2N^{-k}.
\end{equation}
For completeness, use the operator
$(2\pi iD_{v_0}\Psi)^{-1}D_{v_0}$ on the exponential.  The lower
bound for $D_{v_0}\Psi$ and the higher derivative bounds imply that
each fixed directional derivative of $(D_{v_0}\Psi)^{-1}$ is
$O_k(\Lambda^{-1})$.  Applying the adjoint $k$ times to the fixed
smooth bump therefore costs $O_k(\Lambda^{-k})$.  No upper bound on
the first derivative is required.  In particular, these constants
remain uniform as $H$ becomes arbitrarily large: the lower bound and
the higher derivative bounds carry the same factor $H$.

There are at most $C_ZN^2$ retained spatial labels in each ball block.
Thus every row and column of the off-diagonal block matrix has at
most $C_Z\mu^{-8}N^2\leq C_ZN^3$ entries.  Taking $k=4$ in
Equation~\eqref{eq:sparse-kernel-decay} and applying the row and
column sum bound for the $\ell^2$ operator norm shows that its norm
is at most $C\delta^2$.  Together with the diagonal block estimate,
this gives $\|T_\nu\|_{2\to2}^2\leq C\delta^2$.  Applying it to
$h_*\mathbf1_{\Gamma_\nu}$ proves
Equation~\eqref{eq:sparse-joint-mass}.

Since the windows have bounded overlap and
$|\Gamma_\nu|\leq C\delta^2$, H\"older's inequality gives the
precise power needed for Proposition~\ref{prop:discrete}:
\begin{align}
 \delta^2\sum_{\nu,B}(b_\nu^{B,u})^{3/2}
 &\leq\delta^2\sum_\nu\left(\sum_B b_\nu^{B,u}\right)^{3/2}
 \leq C\delta^5\sum_\nu
              \left(\int_{\Gamma_\nu}|h_*|^2\right)^{3/2}
 \notag\\
 &\leq C\delta^5\sum_\nu |\Gamma_\nu|^{1/2}
                     \int_{\Gamma_\nu}|h_*|^3
 \leq C\delta^6.                                      \label{eq:sparse-l3-mass}
\end{align}

\emph{Reconstruction and the discarded spatial tail.}
Let $d^{B,u}$ now denote the finite array in
Equation~\eqref{eq:sparse-analysis} restricted to $|z|\leq Z$.
Apply the synthesis and sampling map to this array, retaining only
$\nu'=\nu_*$, $|z'|\leq1$, and $0\leq j<N^2$.  Write the resulting
array as $c^{B,u}_{z',j}$.  We next justify, uniformly in $B,u,z',j$,
\begin{equation}\label{eq:sparse-samples}
 c^{B,u}_{z',j}
 =F_*\big(a_B+u+(N^2z',j)\big)+O(N^{-3}).
\end{equation}

To do so, write $\widehat d^{B,u}$ for the initial analysis array
with the spatial grid untruncated. Lemma~\ref{lem:packet-frame}
at the initial scale gives
\begin{equation}\label{eq:sparse-full-frame}
 (D\delta)^{-2}\sum_{\nu,z}
       \widehat d^{B,u}_{\nu,z}\chi_\delta^\nu(\xi)
                       e(-N^2z\cdot\xi)=h_{B,u}(\xi)
       \quad\hbox{in }L^2(\mathbb R^2).
\end{equation}
The same lemma gives the exact squared-mass identity
\begin{equation}\label{eq:sparse-full-mass}
 \sum_{\nu,z}|\widehat d^{B,u}_{\nu,z}|^2
  =(D\delta)^2\|h_{B,u}\|_2^2
  =(D\delta)^2\|h_*\|_2^2\leq C\delta^2.
\end{equation}
These statements require only $h\in L^2(Q_0)$, which follows from
the assumed $L^3$ bound and compact support.  In particular, the
possibly very large center modulation does not affect the bound.

A matrix entry from an omitted $|z|>Z$ to a retained output is
\[
 A_{z',j;\nu,z}
  =(D\delta)^{-2}\int\chi_\delta^\nu(\xi)\chi_1^{\nu_*}(\xi)
          e\big(N^2(z'-z)\cdot\xi+j\phi(\xi)\big)\,d\xi.
\]
After the substitution $\xi=\nu+\delta v$, the prefactor is
$D^{-2}$, the amplitude has uniformly bounded smooth derivatives,
and the derivative of the phase in the direction $-z/|z|$ is at least
\[
 N(|z|-1-G)\geq cN(1+|z|).
\]
All phase derivatives of order at least two are bounded uniformly,
because $j\delta^2\leq1$.  The same integration by parts argument
therefore gives
\begin{equation}\label{eq:sparse-tail-entry}
 |A_{z',j;\nu,z}|\leq C_k[N(1+|z|)]^{-k}.
\end{equation}
There are $O(N^2)$ angular labels, and on the spatial mesh
$D^{-1}N^{-1}\mathbb Z^2$,
\[
 \sum_{|z|>Z}(1+|z|)^{-2k}\leq C_kN^2\qquad(k>2).
\]
Consequently each discarded part of a matrix row has squared
$\ell^2$ norm at most $C_kN^{4-2k}$.  By
Equation~\eqref{eq:sparse-full-mass}, its pairing with
$\widehat d^{B,u}$ is at most $C_kN^{1-k}$.  Taking $k=4$ proves
Equation~\eqref{eq:sparse-samples}.  One may justify all infinite
sums by the $L^2$ convergence in
Equation~\eqref{eq:sparse-full-frame}; the tail pairing is also
absolutely convergent by Cauchy--Schwarz.  Thus neither the data nor
the sharp excision boundary has been differentiated.

\emph{Summing samples and filling the balls.}
Apply Proposition~\ref{prop:discrete} to each retained array and sum
over $B$.  Equation~\eqref{eq:sparse-l3-mass} cancels its factor
$N^6$ exactly and gives
\[
 \sum_{B,z',j}|c^{B,u}_{z',j}|^3
 \leq C P(\kappa,\eta)M(s)^{3/2}
                         N^{\gamma_0}\log^2(2+N).
\]
There are $O(N^4)$ terminal spatial labels and $N^2$ times per ball,
so the total number of samples is $O(\mu^{-8}N^6)=O(N^7)$.
The cubed errors from Equation~\eqref{eq:sparse-samples} sum to at
most $CN^7N^{-9}=CN^{-2}$.  Using
$|a+b|^3\leq4(|a|^3+|b|^3)$ and $P,M\geq1$, we obtain
\begin{equation}\label{eq:sparse-cube-sum}
 \sum_{B,z',j}
  |F_*\big(a_B+u+(N^2z',j)\big)|^3
 \leq C P(\kappa,\eta)M(s)^{3/2}
                         N^{\gamma_0}\log^2(2+N),
\end{equation}
uniformly for $u\in\mathcal U$.

Integrating in the common offset $u$ turns the sum for each ball
into an integral over the disjoint half-open cells
\[
 a_B+(N^2z',j)+\mathcal U,
 \qquad |z'|\leq1,\quad 0\leq j<N^2.
\]
These cells cover $B$.  Indeed, for $y\in B$ set $r=y-a_B$; then
\[
 |r'|\leq R\leq N^2/4,\qquad
 N^2/4\leq r_3\leq3N^2/4.
\]
The lower corner $\ell'$ of its horizontal cell in
$D^{-1}\mathbb Z^2$ satisfies
$|\ell'|\leq N^2/4+\sqrt2/D<N^2$ for the fixed large $N$ under
consideration.  Hence $z'=\ell'/N^2$ is a retained label, and
$j=\lfloor r_3\rfloor$ is one of the retained times.  The change
of variables has Jacobian one; the offset domain has fixed volume
$D^{-2}$.  Equation~\eqref{eq:sparse-cube-sum} therefore implies
\[
 \sum_{B\in\mathcal B}\int_B|F_*(y)|^3\,dy
 \leq C P(\kappa,\eta)M(s)^{3/2}
                         N^{\gamma_0}\log^2(2+N).
\]
On $\{|F|>\mu\}$, Equation~\eqref{eq:sparse-excision} gives
$|F_*|>\mu/2$.  Chebyshev's inequality now proves
Equation~\eqref{eq:sparse-estimate}.

All choices made in this proof involve only fixed geometric
constants and fixed orders of differentiation.  Thus $\mu_0$ is
independent of $\kappa,\eta,s$, as asserted.
\end{proof}
